\documentclass[11pt]{article}

\input{Shared.tex}

\parindent0in
\pagestyle{plain}
\thispagestyle{plain}

\usepackage{csquotes}
\usepackage[shortlabels]{enumitem}

\newcommand{\dated}{\today}
\newcommand{\token}[1]{\langle \text{#1} \rangle}

\begin{document}

\textbf{Introduction to the Theory of
Computation}\hfill\textbf{\myname}\\[0.01in]
\textbf{Chapter 8: Space Complexity}\hfill\textbf{\dated}\\
\smallskip\hrule\bigskip

\begin{problem}{8.10}
The Japanese game \textit{go-moku} is played by two players, ``X'' and ``O'' on a \mbox{$19 \times 19$} grid. Players take turns placing markers, and the first player to achieve five of her markers consecutively in a row, column, or diagonal is the winner. Consider this game generalized to an $n \times n$ board. Let
\[
GM = \{\langle B \rangle \ | \ B \text{ is a position in generalized go-moku, where player ``X'' has a winning strategy}\}.
\]
By a $position$ we mean a board with markers placed on it, such as may occur in the middle of a play of the game, together with an indication of which player moves next. Show that $GM \in PSPACE$.
\end{problem}

\begin{proof}
\end{proof}

\end{document}